\section*{Guion principal: 29 minutos}

\NarrativeEntry{1}{Portada y apertura}{0:30}{0:30}{
Buenos días. Esta defensa presenta un estudio sobre reconstrucción de canales EEG faltantes mediante procesamiento de señales en grafos y regularización temporal. La pregunta que organiza toda la exposición es simple de formular, pero exige evidencia cuidadosa: ¿cuándo la continuidad temporal agrega información útil a una referencia espacial madura como MNE-Python? Voy a responderla desde el problema, el diseño experimental, los resultados y, finalmente, las condiciones bajo las cuales conviene elegir cada método.
}{Mirar al jurado. Pausa breve después de la pregunta.}{Tesis final, portada e Introducción.}

\NarrativeEntry{2}{El problema real}{1:10}{1:40}{
Un canal EEG puede perderse por mal contacto, desplazamiento o artefactos persistentes. Descartar todo el registro reduce la información disponible; reconstruir el canal evita ese descarte, pero introduce una decisión metodológica antes de cualquier análisis posterior. Por eso no considero la interpolación como un simple relleno. Una reconstrucción modifica la topografía, la forma temporal y el contenido espectral que luego alimentan otras etapas. Para medir ese efecto con verdad de terreno, el estudio parte de señales completas y oculta canales de manera artificial. Así conocemos exactamente qué señal debió recuperarse y podemos comparar métodos bajo la misma máscara. Este control fortalece la validez interna, aunque no representa por completo un electrodo realmente dañado; esa frontera aparecerá nuevamente en las limitaciones.
}{Señalar la cadena de impacto de izquierda a derecha.}{Tesis, Introducción y protocolo de ocultamiento artificial.}

\NarrativeEntry{3}{Por qué no es trivial}{1:05}{2:45}{
La reconstrucción es difícil por cuatro razones. Primero, la información espacial depende de qué electrodos permanecen observados: una falla agrupada puede eliminar varios vecinos útiles al mismo tiempo. Segundo, el EEG tiene continuidad temporal, pero regularizarla en exceso puede suavizar componentes de interés. Tercero, las métricas no son equivalentes. Reducir el error absoluto no garantiza preservar mejor el espectro ni ejecutar más rápido. Cuarto, el método adecuado puede cambiar con la severidad y el patrón de pérdida. El problema, entonces, no es encontrar una técnica que gane siempre, sino caracterizar un régimen de decisión: cuándo la estructura espacio--temporal compensa su costo y cuándo una referencia espacial consolidada sigue siendo suficiente.
}{Marcar una pausa entre cada una de las cuatro razones.}{Tesis, caps. 1, 2 y 5.}

\NarrativeEntry{4}{Brecha}{1:00}{3:45}{
La brecha metodológica tiene tres componentes. Una comparación contra un spline propio puede ser útil para desarrollo, pero no sustituye una implementación comunitaria fuerte. Una evaluación basada solo en MAE oculta tensiones con forma temporal, fidelidad espectral y tiempo. Y una prueba con pérdidas aleatorias leves no cubre fallas localizadas. La tesis responde combinando un comparador fuerte, MNE-Python; un portafolio de métricas; y varios regímenes de pérdida. Además, separa la exploración de la evaluación final. Esa separación es crucial: la amplitud de dieciocho métodos, nueve grafos y ciento veinte escenarios orienta decisiones, pero las cifras centrales provienen únicamente del protocolo congelado TRSS fijo frente a MNE-Python.
}{Enfatizar “protocolo congelado”.}{Tesis, Introducción y Metodología.}

\NarrativeEntry{5}{Preguntas y objetivo}{1:05}{4:50}{
La primera pregunta busca las condiciones en que la regularización espacio--temporal mejora la reconstrucción frente a métodos espaciales de referencia. La segunda pregunta examina el compromiso práctico entre precisión, fidelidad espectral y tiempo de cómputo al elegir entre TRSS y MNE-Python. El objetivo general es comparar esas dimensiones bajo un marco riguroso y reproducible. No planteo una superioridad universal ni una validación clínica. La tesis busca una respuesta condicional y trazable: identificar dónde aparece la ventaja, cuantificar sus contrapartidas y explicitar qué evidencia sería necesaria para extender el alcance.
}{Leer P1 y P2 casi literalmente; luego resumir el objetivo con voz propia.}{Tesis, Introducción, preguntas P1 y P2.}

\NarrativeEntry{6}{EEG como señal sobre un grafo}{1:15}{6:05}{
El procesamiento de señales en grafos representa cada electrodo como un nodo y la relación entre electrodos como aristas ponderadas. Una muestra EEG es entonces una señal definida sobre esos nodos. El Laplaciano del grafo cuantifica variación espacial: si dos nodos conectados toman valores muy distintos, la energía de suavidad aumenta. Cuando falta un canal, este formalismo permite reconstruirlo buscando una señal compatible con los canales observados y, a la vez, regular respecto de la topología. La ventaja conceptual es que la noción de vecindad queda explícita y puede inspeccionarse. En la comparación final se privilegian grafos geométricos deterministas, construidos desde coordenadas de electrodos, para no usar la señal evaluada en la definición de sus propias relaciones.
}{No explicar toda la teoría espectral; señalar nodos, aristas y Laplaciano.}{Tesis, cap. 2 y figura de conceptos GSP.}

\NarrativeEntry{7}{MNE frente a TRSS}{1:25}{7:30}{
Los dos métodos finales usan información diferente. MNE-Python interpola mediante splines esféricos y añade decisiones de ingeniería como normalización, restricciones, pseudoinversa y manejo robusto de casos numéricos. Por eso es una referencia más exigente que un spline académico simplificado. TRSS resuelve un problema conjunto sobre canales y tiempo. El primer término obliga a respetar las observaciones; el segundo, controlado por alfa, penaliza variación espacial sobre el grafo; y el tercero, controlado por beta, penaliza cambios abruptos entre muestras. Si beta es cero, queda una regularización espacial independiente por instante. La hipótesis operativa no es que temporalidad siempre sea mejor, sino que puede aportar cuando la geometría observada deja poca información local. La misma continuidad puede transformarse en desventaja si suaviza de forma inapropiada o si su costo domina.
}{Señalar los tres términos del diagrama, no desarrollar la solución algebraica.}{Tesis, cap. 2, formulación TRSS y mecanismo de MNE.}

\NarrativeEntry{8}{Decisiones de alcance}{1:00}{8:30}{
Antes de mostrar experimentos, explicito cuatro decisiones. MNE-Python es la línea base principal para evitar comparadores debilitados. Los grafos finales son geométricos y reproducibles; los aprendidos quedan como trabajo futuro. No se incluyeron modelos entrenables porque no había un protocolo externo de entrenamiento que permitiera compararlos sin abrir otro problema experimental. Finalmente, la lectura principal usa TRSS fijo, con parámetros congelados antes de la evaluación. Estas decisiones restringen el alcance, pero también protegen la interpretación. No significan que las alternativas excluidas sean incapaces; significan que esta tesis no aporta evidencia suficiente para compararlas de manera justa.
}{Mantener tono defensivo pero no apologético.}{Tesis, Metodología y Discusión.}

\NarrativeEntry{9}{Estrategia experimental}{1:25}{9:55}{
El flujo experimental tiene varias capas. Primero se estandarizan y segmentan tres bases públicas. Después se generan máscaras de pérdida con semillas controladas. Luego se construyen grafos, se ejecutan métodos y se calculan métricas. La fase exploratoria recorre dieciocho implementaciones y nueve grafos; la fase intermedia reduce candidatos por estabilidad, costo y coherencia; después se optimizan parámetros y se congelan variantes. Solo entonces se ejecuta la comparación final pareada. Esta secuencia evita confundir búsqueda con confirmación. También permite que cada resultado conserve un identificador, la máscara utilizada, parámetros, tiempos y señales reconstruidas. La trazabilidad no es un apéndice administrativo: es lo que permite regresar desde una cifra de la defensa hasta la configuración que la originó.
}{Recorrer visualmente el flujo y detenerse en el congelamiento.}{Tesis, figura de flujo metodológico.}

\NarrativeEntry{10}{Datos}{1:05}{11:00}{
El estudio usa tres conjuntos públicos con características complementarias. PhysioNet aporta 64 canales, 160 hertz y 109 sujetos en reposo. BCI Competition IV 2a aporta 22 canales, 250 hertz y nueve sujetos en imaginación motora. MNE Sample aporta 60 canales, 600 hertz y un sujeto en un paradigma auditivo. Esta diversidad prueba montajes, frecuencias y paradigmas diferentes. Sin embargo, no debe interpretarse como una muestra representativa de toda población EEG. Especialmente MNE Sample contribuye diversidad técnica, no evidencia poblacional. La conclusión válida es que el comportamiento se observó bajo varios entornos públicos controlados; extenderlo a otros equipos, montajes o poblaciones exige validación adicional.
}{Contrastar explícitamente diversidad técnica con generalización poblacional.}{Tesis, Metodología, tabla de conjuntos de datos.}

\NarrativeEntry{11}{Cómo se simula la pérdida}{1:05}{12:05}{
Se estudian dos patrones. En el aleatorio, los canales faltantes se distribuyen; en el contiguo, se elige un canal semilla y se agregan vecinos cercanos, aproximando una falla localizada. La fase amplia cruza tres datasets, dos modos, cuatro severidades y cinco semillas: ciento veinte escenarios. Dentro de cada escenario, todos los métodos reciben exactamente la misma señal y la misma máscara. Esto permite comparaciones pareadas y evita que un método se beneficie de una pérdida más sencilla. La ventaja del ocultamiento artificial es conocer la verdad; su límite es que un canal real puede tener ruido residual, deriva o artefactos, no solo ausencia completa.
}{Explicar primero el dibujo mental de las máscaras y luego la aritmética 3 por 2 por 4 por 5.}{Tesis, Metodología, generación de máscaras.}

\NarrativeEntry{12}{Evaluación multi-métrica}{1:05}{13:10}{
La evaluación separa dimensiones. MAE, RMSE y NRMSE miden error de amplitud; SNR la relación entre señal y error; DTW, correlación y R cuadrado capturan aspectos de forma temporal; LSD, coherencia y PSD examinan el dominio espectral; y el tiempo caracteriza costo. Cada métrica tiene una dirección favorable definida antes del análisis. No construyo un puntaje único, porque eso obligaría a imponer pesos subjetivos entre objetivos distintos. En su lugar, presento un portafolio y pregunto si la misma conclusión se sostiene en cada eje. Esta elección permite reconocer que un método puede mejorar la amplitud y, simultáneamente, perder en fidelidad espectral o velocidad.
}{Nombrar familias, no recitar fórmulas.}{Tesis, Metodología, sección Métricas.}

\NarrativeEntry{13}{Arquitectura y trazabilidad}{1:00}{14:10}{
La arquitectura implementa esa lógica con configuración inmutable, carga de datos, construcción del grafo, ejecución del interpolador, cálculo de métricas y persistencia de artefactos. Cada caso conserva dataset, patrón, severidad, semilla, máscara, parámetros, tiempos y reconstrucción. Un identificador SHA-256 permite detectar cambios y vincular filas, figuras y configuraciones. También se validan dimensionalidad, canales observados y faltantes, coherencia temporal y correspondencia entre máscara y semilla. El objetivo es que una cifra no dependa de memoria manual ni de una captura aislada: debe poder localizarse, regenerarse y auditarse.
}{Señalar el recorrido entrada--resultado; no entrar a detalles de código.}{Tesis, Metodología y figura de arquitectura.}

\NarrativeEntry{14}{Reducción experimental}{1:10}{15:20}{
La amplitud inicial no equivale a la comparación final. La exploración consideró dieciocho implementaciones agrupadas en métodos geométricos, estadísticos, GSP y espacio--temporales, además de nueve grafos. Luego se redujeron candidatos por precisión, estabilidad numérica, costo y reproducibilidad. Algunas opciones aprendidas se descartaron por costo y riesgo de circularidad; los métodos geométricos simples sirvieron como referencia exploratoria; TRSS avanzó por su potencial espacio--temporal. La fase de optimización ajustó familias, pero la evaluación final congeló TRSS fijo y MNE-Python. La variante calibrada y el oráculo se conservan solo para análisis auxiliar. Por eso no mezclo cifras exploratorias con las que responden P1 y P2.
}{Usar la caída de 18 métodos a 2 como argumento de rigor, no como ranking total.}{Tesis, tablas de fases y decisiones de reducción.}

\NarrativeEntry{15}{Protocolo final}{1:25}{16:45}{
El protocolo final contiene cien estratos definidos por régimen, patrón y severidad. En cada estrato hay tres semillas, para un total de trescientas diferencias pareadas TRSS--MNE. Dentro de cada par se comparten señal, máscara y preprocesamiento; la diferencia se calcula en el mismo caso. La unidad observacional es esa diferencia por estrato y semilla, con $n=300$; la unidad de agrupación es el estrato, con $n=100$. Para resumir estabilidad se usa un bootstrap jerárquico descriptivo de mil remuestreos que respeta ese agrupamiento. Las semillas no son sujetos nuevos. Los intervalos describen la variabilidad del protocolo; no son pruebas de hipótesis ni justifican inferencia poblacional externa. Esta precisión evita inflar el tamaño efectivo de la evidencia.
}{Enfatizar 100, 3 y 300; luego aclarar qué no significa el intervalo.}{Tesis, Protocolo de Validación Comparativa.}

\NarrativeEntry{16}{Resultado principal en MAE}{1:30}{18:15}{
El resultado central es una mejora mediana de 12,4 por ciento en MAE para TRSS fijo frente a MNE-Python. El intervalo descriptivo al 95 por ciento va de 7,8 a 17,4 por ciento y TRSS gana en 72 por ciento de los trescientos casos. Esta es la evidencia más directa a favor de la regularización temporal bajo el protocolo congelado. Pero la cifra no significa que cada escenario mejore ni que la diferencia se transfiera automáticamente a otro montaje. Resume el conjunto de casos diseñados aquí. La tasa de victoria complementa la mediana: muestra que la tendencia no depende de un pequeño número de mejoras extremas, aunque todavía existe una proporción relevante de empates o victorias de MNE.
}{Pausa después de 12,4 por ciento. Señalar intervalo y tasa de victoria.}{Tesis, tabla principal robusta y figura de intervalos.}

\NarrativeEntry{17}{Portafolio de métricas}{1:30}{19:45}{
Cuando ampliamos la mirada, TRSS conserva ventajas agregadas en NRMSE, DTW, SNR, correlación y R cuadrado. Por ejemplo, la mejora mediana es 13,9 por ciento en NRMSE y 9,8 por ciento en DTW. Sin embargo, LSD muestra una mejora mediana de menos 2,3 por ciento y una tasa de victoria TRSS de 40,7 por ciento: en esa métrica global, MNE es mejor. El tiempo también favorece claramente a MNE. Esta tensión es parte del resultado, no una anomalía que deba ocultarse. La respuesta correcta a P2 no es declarar un ganador, sino reconocer una frontera: TRSS mejora varios aspectos de amplitud y forma temporal, mientras MNE preserva mejor el LSD agregado y ejecuta mucho más rápido.
}{Mostrar primero el grupo favorable y después detenerse en LSD y tiempo.}{Tesis, tabla y figura del portafolio de métricas.}

\NarrativeEntry{18}{Heterogeneidad por escenario}{1:30}{21:15}{
La ventaja no se distribuye uniformemente. En pérdidas cercanas o agrupadas con 30 y 40 por ciento de canales faltantes, la mejora mediana en MAE supera 40 por ciento y la tasa de victoria alcanza 80 u 86,7 por ciento. En pérdidas periféricas severas también aparece una ventaja importante. En cambio, con dos canales cercanos faltantes, la mediana puede ser cero y la tasa de victoria queda próxima al empate. La interpretación mecánica es coherente: cuando desaparecen vecinos informativos, la continuidad temporal aporta una restricción adicional; cuando la geometría espacial conserva suficiente información, el costo y el suavizado extra pueden no aportar. Estos estratos ilustran condiciones frontera, pero no estiman con qué frecuencia aparecerán en una práctica externa.
}{Comparar un extremo favorable y el caso de empate; evitar leer toda la matriz.}{Tesis, tabla de escenarios seleccionados y mapa de calor.}

\NarrativeEntry{19}{Forma de señal}{1:10}{22:25}{
Las series temporales representativas ayudan a entender el mecanismo, pero no reemplazan el agregado. Se observan casos donde TRSS sigue mejor la forma original, otros donde MNE es suficiente y algunos donde ambos fallan de manera distinta. Las curvas permiten detectar suavizado, desfase o desviaciones locales que un número resume parcialmente. Por eso se usan como ejemplos auditables y no como prueba de superioridad. La conclusión cuantitativa sigue descansando en los trescientos pares y en el análisis por estratos. El valor de esta lámina es conectar la estadística con una señal concreta y recordar qué significa, en términos físicos, reducir el error.
}{Señalar original y reconstrucciones; no inferir desde un único canal.}{Tesis, Resultados, series temporales representativas.}

\NarrativeEntry{20}{Tensión espectral}{1:10}{23:35}{
El análisis de densidad espectral muestra por qué era necesario un portafolio. Un método puede aproximar mejor la señal muestra a muestra y, aun así, alterar la distribución de energía entre frecuencias. El LSD agregado favorece a MNE: TRSS tiene una mejora mediana negativa de 2,3 por ciento y gana solo en 40,7 por ciento de los casos. Esto no invalida el resultado en MAE; indica que los objetivos no son intercambiables. Si la prioridad de una aplicación es preservar el espectro global, MNE tiene una ventaja bajo este benchmark. Si importa la reconstrucción temporal en una pérdida localizada y el costo es aceptable, TRSS puede ser preferible.
}{Conectar explícitamente PSD cualitativa con LSD cuantitativo.}{Tesis, Resultados, PSD y portafolio de métricas.}

\NarrativeEntry{21}{Costo computacional}{1:05}{24:40}{
La diferencia de costo es clara. MNE-Python presenta una mediana de 0,0090 segundos por caso; TRSS fijo, 0,1214 segundos. En esta implementación y hardware, TRSS es aproximadamente un orden de magnitud más lento. Los tiempos no son constantes universales: dependen del solver, tamaño, hardware e ingeniería. Aun así, la dirección de la diferencia es suficientemente estable para formar parte de la decisión. La variante calibrada no elimina el costo de búsqueda de hiperparámetros y no es la lectura principal. Por tanto, un sistema de alto volumen o con restricción de latencia debe privilegiar MNE, salvo que la dificultad del patrón de pérdida justifique el costo adicional.
}{Dar cifras con cuatro decimales y aclarar dependencia de implementación.}{Tesis, tabla de tiempo y complejidad.}

\NarrativeEntry{22}{Frontera práctica}{1:05}{25:45}{
La evidencia se traduce en una regla condicional. MNE es la opción por defecto cuando faltan pocos canales, la geometría conserva vecinos útiles, se prioriza LSD global o existe una restricción fuerte de tiempo. TRSS merece consideración cuando la pérdida es espacialmente difícil, faltan vecinos cercanos o periféricos, el análisis temporal es importante y el procesamiento fuera de línea admite mayor costo. Esta regla no automatiza una decisión clínica ni reemplaza validación posterior. Es un mapa de diseño basado en el benchmark: permite elegir de manera explícita qué objetivo se prioriza, en vez de asumir que una sola métrica decide por todos.
}{Recorrer los dos lados del mapa y cerrar en “regla condicional”.}{Tesis, Discusión y mapa de decisión.}

\NarrativeEntry{23}{Aportes}{1:00}{26:45}{
El trabajo aporta cuatro elementos. Primero, una caracterización condicional de cuándo la regularización temporal agrega valor. Segundo, una metodología que separa exploración, reducción, congelamiento y evaluación final. Tercero, una comparación pareada y multi-métrica contra una referencia comunitaria fuerte. Cuarto, una cadena reproducible desde la configuración y la máscara hasta la tabla, la figura y la afirmación. La novedad no descansa en inventar el procesamiento de grafos ni los splines, sino en integrar TRSS, un benchmark exigente y una lectura de frontera respaldada por trazabilidad.
}{Enumerar con la mano; no sobredimensionar la novedad algorítmica.}{Tesis, Conclusiones, sección Aportes.}

\NarrativeEntry{24}{Limitaciones}{1:15}{28:00}{
Las limitaciones definen qué no concluyo. El ocultamiento artificial entrega verdad de terreno, pero no reproduce ruido residual ni deterioro real. Las medidas repetidas permiten comparación pareada, pero no equivalen a sujetos independientes ni sostienen inferencia poblacional. Las métricas son proxies de reconstrucción; falta evaluar impacto en clasificación, conectividad o lectura experta. Los casos visuales no son una evaluación ciega. Los tiempos pertenecen a una implementación y hardware reportados. Tampoco se compararon modelos entrenables bajo un protocolo externo. El siguiente paso debe incluir daño real, tareas posteriores, evaluación ciega, otros montajes y una versión acelerada de TRSS. Presentar estos límites no debilita la tesis: evita que una evidencia controlada se transforme en una promesa que no fue ensayada.
}{Reducir velocidad y mirar al jurado; esta lámina prepara preguntas difíciles.}{Tesis, Discusión y Conclusiones.}

\NarrativeEntry{25}{Conclusión}{1:00}{29:00}{
Concluyo con dos respuestas. Para P1, TRSS fijo mejora la reconstrucción en el agregado y aporta especialmente cuando la pérdida espacial elimina vecinos informativos; no reemplaza universalmente a MNE. Para P2, la ganancia en amplitud y forma temporal se intercambia por mayor costo y una desventaja en LSD global. La decisión depende del régimen y del objetivo posterior. El aporte central es haber identificado esa frontera con un protocolo pareado, congelado y trazable, mostrando tanto la ventaja como sus límites. Muchas gracias; quedo atento a sus preguntas.
}{Pausa entre P1 y P2. Terminar mirando al jurado, no a la pantalla.}{Tesis, Conclusiones y respuestas a P1/P2.}
